% Generated draft: requires mathematical and source review.
\documentclass[11pt]{article}
\usepackage[margin=25mm]{geometry}
\usepackage{fontspec}
\setmainfont{Latin Modern Roman}
\usepackage{amsmath,amssymb,mathtools,unicode-math}
\setmathfont{Latin Modern Math}
\usepackage{enumitem,needspace,xurl,hyperref}
\setlength{\parindent}{0pt}
\setlength{\parskip}{5pt}
\setlength{\emergencystretch}{3em}
\newcommand{\N}{\mathbb{N}}
\newcommand{\Z}{\mathbb{Z}}
\newcommand{\Q}{\mathbb{Q}}
\newcommand{\R}{\mathbb{R}}
\newcommand{\C}{\mathbb{C}}
\newcommand{\F}{\mathbb{F}}
\newcommand{\A}{\mathbb{A}}
\newcommand{\PP}{\mathbb{P}}
\newcommand{\E}{\mathbb{E}}
\newcommand{\eps}{\varepsilon}
\newcommand{\dd}{\,\mathrm{d}}
\begin{document}
\section*{UnsolvedMath ID: 30001264}
\textbf{Source problem code:} catalogue entry
\subsection{Definitions and mathematical statement}
Consider a central extension
\[
0\longrightarrow\Z\xrightarrow{\iota}E\xrightarrow{\pi}G\longrightarrow1,
\qquad \ker\pi=\iota(\Z)\subseteq Z(E).
\]
For the abstract-group interpretation, give \(\Z\) the trivial \(G\)-action. Choose a set-theoretic section \(s:G\to E\) with \(s(1)=1\). Its defining cocycle \(c:G\times G\to\Z\) is determined by
\[
s(g)s(h)=\iota(c(g,h))s(gh).
\]
It satisfies
\[
c(g,h)+c(gh,k)=c(h,k)+c(g,hk)
\qquad(g,h,k\in G),
\]
and \(c(1,g)=c(g,1)=0\). Changing the section replaces \(c\) by \(c+\delta b\), where
\[
(\delta b)(g,h)=b(g)+b(h)-b(gh),
\]
so the extension determines a class \([c]\in H^2(G;\Z)\).

The source does not specify the metric framework for the word ``undistorted.'' Under the usual interpretation when \(E\) is a finitely generated discrete group, choose a finite generating set \(S\) and let \(|\cdot|_S\) be its word length. The central subgroup is undistorted if its inclusion is a quasi-isometric embedding: there exist \(A\geq1\) and \(B\geq0\) such that
\[
A^{-1}|n|-B\leq |\iota(n)|_S\leq A|n|+B
\qquad\text{for every }n\in\Z.
\]
This interpretation is conditional; finite generation and discreteness are not stated in the supplied question.

\textbf{Source question.} For a central \(\Z\)-extension whose central subgroup is undistorted, determine what restrictions this imposes on the cohomology class of its defining two-cocycle, interpreted as \([c]\in H^2(G;\Z)\) in the abstract-group setting above.

\textbf{Unresolved qualifications.} The supplied formulation does not specify a class of groups, a metric or length function outside the finitely generated setting, or a cohomology theory for topological groups. It also does not state a particular proposed restriction on \([c]\). Thus it supplies a classification question rather than a definite conjectural implication. These omissions prevent a unique fully specified general formulation. The supplied literature assessment concerns connected Lie groups, but does not provide the hypotheses or conclusion of the claimed restricted result.
\subsection{Short English statement}
The question asks how an undistorted central copy of \(\Z\) constrains the cohomology class defining the extension. The group category, metric assumptions, and intended cohomological restriction remain unspecified. The supplied record labels the problem partially solved and reports a complete answer for connected Lie groups; this is dated curation requiring independent verification.
\subsection{Sources}
\begin{itemize}
\item Local UnsolvedMath registry, \texttt{data/unsolvedmath/problems.json}, permanent ID \(30001264\), source problem code \texttt{catalogue entry}, title: ``Cohomology Classes of Undistorted Central Extensions.'' The supplied record attributes the problem to the Oberwolfach Reports collection, with proposed year 2009 and no named proposer.
\item Supplied report citation: \emph{Manifold Perspectives} (2009), Oberwolfach Reports, pp.~1487--1546, DOI 10.4171/owr/2009/27. \url{https://doi.org/10.4171/owr/2009/27}. The supplied background describes these reports as documenting mathematical workshops and their open-problem discussions; no exact problem page is supplied.
\item Supplied literature source, labeled ``bounded cohomology of Lie groups'': \url{https://arxiv.org/abs/0904.3841}.
\item Supplied literature source, labeled ``distortion and bounded cohomology'': \url{https://chatterj.perso.math.cnrs.fr/papers/mfoMay09.pdf}.
\item Supplied literature source, labeled ``original OWR report'': \url{https://doi.org/10.4171/owr/2009/27}.
\item Literature assessment dated exactly \texttt{2026-08-21}: the supplied status is \texttt{partially\_solved}, and the supplied assessment says that the cited work gives a complete answer for connected Lie groups while the extracted question concerns general groups. This is supplied dated curation derived from the OWR record, not an independently verified literature assessment. The supplied background expressly requests independent verification of time-sensitive status and cited claims.
\end{itemize}
\end{document}
